\section{Evaluation Design}\label{sec:evaluation}
\begin{table}[H]\centering\small
\caption{Research questions and evidence boundaries.}
\begin{tabularx}{\linewidth}{@{}p{.06\linewidth}p{.35\linewidth}X@{}}\toprule
RQ & Question & Evidence\\\midrule
RQ1 & When does ordering reduce hierarchical packet cost? & Frozen synthetic workloads with disjoint training, validation and held-out epochs\\
RQ2 & What does practical compilation trade for simplicity, and is semantic output invariant? & Archived source re-encodings and local authorization checks\\
RQ3 & Is a beyond-pairwise algorithmic benefit demonstrated? & Re-analysis of the unchanged ablation, with corrected signs, subsets and percentage units\\
RQ4 & When can delivery savings repay extra work? & Explicit conditional break-even scenarios\\
RQ5 & What changes at equal diagnostic service? & Corrected tree and indexed-ledger queries against an independent oracle; separate resource ledgers\\\bottomrule
\end{tabularx}
\end{table}
The v5/v7 records and the v8 ablation are archived measurements. The v9
query study is a new retrospective execution of corrected service code over
the same source records, not new checker trials or a preregistered prospective
study. Its input files are digest-bound. Host elapsed timings are descriptive;
network arrival values are explicitly modeled. Current result macros, query
rows and signed ablation prose are generated by an independent arithmetic audit.
